\documentclass[a4paper,12pt]{article}

\usepackage{iftex}
\ifPDFTeX
  \usepackage[T2A]{fontenc}
  \usepackage[utf8]{inputenc}
  \usepackage[russian]{babel}
  \usepackage{helvet}
  \usepackage{microtype}
\else
  \usepackage{fontspec}
  \setmainfont{DejaVu Sans}[AutoFakeSlant=0.18]
  \setsansfont{DejaVu Sans}[AutoFakeSlant=0.18]
  \renewcommand{\today}{15 июля 2026 г.}
\fi
\renewcommand{\familydefault}{\sfdefault}
\renewcommand{\contentsname}{Содержание}
\usepackage{amsmath,amssymb,amsthm,mathtools}
\usepackage{icomma}
\usepackage[a4paper,left=18mm,right=18mm,top=18mm,bottom=20mm]{geometry}
\usepackage{tikz}
\usepackage{tkz-euclide}
\usepackage{pgfplots}
\pgfplotsset{compat=1.18}
\usepackage{tabularx,booktabs,longtable}
\usepackage{enumitem}
\usepackage{hyperref}
\usepackage{parskip}
\usepackage{fancyhdr}
\usepackage{setspace}
\usepackage{xcolor}
\usepackage{array}

\definecolor{accent}{HTML}{7B3142}
\definecolor{darkaccent}{HTML}{4E2430}
\definecolor{textgray}{HTML}{30343B}
\definecolor{softgray}{HTML}{6E7279}
\definecolor{paper}{HTML}{F5EFE5}

\hypersetup{colorlinks=true,linkcolor=darkaccent,urlcolor=accent}
\onehalfspacing
\setlength{\parindent}{0pt}
\setlength{\parskip}{5pt}
\setlength{\emergencystretch}{3em}
\setlist[itemize]{leftmargin=6mm,itemsep=2pt,topsep=3pt}
\setlist[enumerate]{leftmargin=8mm,itemsep=4pt,topsep=4pt}
\renewcommand{\arraystretch}{1.35}

\pagestyle{fancy}
\fancyhf{}
\lhead{\textcolor{textgray}{Показательные неравенства}}
\rhead{\textcolor{textgray}{15.07.26}}
\cfoot{\textcolor{softgray}{\thepage}}
\setlength{\headheight}{15pt}

\newcommand{\key}[1]{\textcolor{accent}{\textbf{#1}}}
\newcommand{\R}{\mathbb{R}}
\newcommand{\Z}{\mathbb{Z}}

\begin{document}

\begin{titlepage}
  \centering
  \vspace*{24mm}
  {\Large\textcolor{accent}{ЧЕК-ЛИСТ}\par}
  \vspace{8mm}
  {\fontsize{30}{36}\selectfont\bfseries Показательные\\[2mm] неравенства\par}
  \vspace{8mm}
  {\large Общее основание, замена переменной\\и логарифмический переход\par}
  \vfill
  {\large Лёвин Артём Александрович\\эксклюзивно для Тимофея Васильченко\par}
  \vspace{8mm}
  {\large \today\par}
  \vfill
  \begin{tikzpicture}[remember picture,overlay]
    \draw[accent,line width=2pt]
      ([xshift=22mm,yshift=24mm]current page.south west)
      .. controls +(38mm,22mm) and +(-35mm,-10mm) ..
      ([xshift=-22mm,yshift=35mm]current page.south east);
    \draw[darkaccent,line width=.8pt]
      ([xshift=32mm,yshift=18mm]current page.south west)
      .. controls +(45mm,28mm) and +(-48mm,-14mm) ..
      ([xshift=-30mm,yshift=28mm]current page.south east);
  \end{tikzpicture}
\end{titlepage}

\tableofcontents
\thispagestyle{fancy}
\newpage

\section{Что нужно узнать сразу}

\key{Показательное неравенство} содержит переменную в показателе степени, например
\[
2^x\ge 16,\qquad
\left(\frac13\right)^{2x-1}<\frac1{27}.
\]

Для показательной функции требуется
\[
a>0,\qquad a\ne1.
\]
При этих условиях значение $a^x$ определено и строго положительно для любого $x\in\R$.

\subsection{Монотонность определяет знак}

\begin{table}[ht]
\centering
\caption{Равносильные переходы для одинакового основания}
\begin{tabularx}{\linewidth}{>{\bfseries}p{35mm} X X}
\toprule
Основание & Поведение функции & Переход \\
\midrule
$a>1$ & $a^x$ возрастает & $a^{f(x)}\gtreqless a^{g(x)}\iff f(x)\gtreqless g(x)$ \\
$0<a<1$ & $a^x$ убывает & $a^{f(x)}\gtreqless a^{g(x)}\iff f(x)\lesseqgtr g(x)$ \\
\bottomrule
\end{tabularx}
\end{table}

\key{Главная проверка перед сравнением показателей:} основания должны быть одинаковыми, положительными и отличными от единицы. Слагаемые и лишние множители рядом со степенью сначала устраняют пре\-об\-ра\-зо\-ва\-ни\-я\-ми.

\subsection{Степенная запись}

\[
\sqrt[n]{a^m}=a^{m/n},\qquad
a^m a^n=a^{m+n},\qquad
\frac{a^m}{a^n}=a^{m-n},\qquad
(a^m)^n=a^{mn}.
\]

Полезные представления:
\[
4=2^2,\quad 8=2^3,\quad \sqrt[3]{16}=2^{4/3},
\quad 0{,}2=\frac15=5^{-1},\quad 0{,}25=\frac14=2^{-2}.
\]

\section{Метод 1. Приведение к общему основанию}

\subsection{Простейшие случаи}

\[
2^x\ge2^1\iff x\ge1.
\]
Основание $2>1$, поэтому знак сохраняется.

\[
\left(\frac15\right)^x\ge\left(\frac15\right)^1\iff x\le1.
\]
Основание лежит между нулём и единицей, поэтому знак меняется.

\subsection{Пример с составными числами и корнем}

Решим
\[
4^x\cdot8<\sqrt[3]{16}.
\]
Приведём всё к основанию $2$:
\[
2^{2x}\cdot2^3<2^{4/3}
\iff
2^{2x+3}<2^{4/3}.
\]
Так как $2>1$,
\[
2x+3<\frac43
\iff 2x<-\frac53
\iff x<-\frac56.
\]

\subsection{Пример с корнями в обеих частях}

\[
\sqrt2\cdot8^x\ge\frac4{\sqrt[5]{16}}.
\]
Перепишем каждый объект как степень двойки:
\[
2^{1/2}\cdot2^{3x}\ge\frac{2^2}{2^{4/5}}
\iff
2^{3x+1/2}\ge2^{6/5}.
\]
Следовательно,
\[
3x+\frac12\ge\frac65
\iff 3x\ge\frac7{10}
\iff x\ge\frac7{30}.
\]

\subsection{Дробь в показателе}

Рассмотрим
\[
\left(\frac15\right)^{\frac{2x+3}{x-2}}>5.
\]
Сначала фиксируем ограничение $x\ne2$ и заменяем основание:
\[
5^{-\frac{2x+3}{x-2}}>5^1
\iff
-\frac{2x+3}{x-2}>1.
\]
После нормализации:
\[
\frac{3x+1}{x-2}<0.
\]
Метод интервалов даёт
\[
\boxed{x\in\left(-\frac13;2\right)}.
\]

\section{Метод 2. Замена \texorpdfstring{\(t=a^x\)}{t равно a в степени x}}

\subsection{Когда нужна замена}

Если в неравенстве повторяются $a^x$, $a^{2x}$, $a^{3x}$ или дроби с такими выражениями, вводят
\[
t=a^x,\qquad t>0.
\]
Ограничение $t>0$ обязательно: показательная функция не принимает нулевых и отрицательных значений.

\subsection{Квадратное неравенство после замены}

\[
9^x-10\cdot3^x+9\le0.
\]
Положим $t=3^x>0$. Тогда $9^x=(3^x)^2=t^2$:
\[
t^2-10t+9\le0
\iff
(t-1)(t-9)\le0.
\]
Отсюда $1\le t\le9$. Возвращаемся к $x$:
\[
1\le3^x\le9
\iff
3^0\le3^x\le3^2
\iff
\boxed{x\in[0;2]}.
\]

\subsection{Полный пример экзаменационного уровня}

Решим
\[
\frac7{9^x-2}\ge\frac2{3^x-1}.
\]

\key{1. Ограничения.}
\[
9^x\ne2,\qquad 3^x\ne1.
\]

\key{2. Замена.} Пусть $t=3^x>0$. Тогда $9^x=t^2$:
\[
\frac7{t^2-2}\ge\frac2{t-1}.
\]

\key{3. Нормализация.}
\[
\frac{7(t-1)-2(t^2-2)}{(t^2-2)(t-1)}\ge0,
\]
\[
\frac{-2t^2+7t-3}{(t^2-2)(t-1)}\ge0
\iff
\frac{2t^2-7t+3}{(t^2-2)(t-1)}\le0.
\]

\key{4. Критические точки.}
\[
2t^2-7t+3=0
\iff t=\frac12\ \text{или}\ t=3.
\]
Знаменатель обращается в ноль при $t=1$ и $t=\pm\sqrt2$. С учётом $t>0$ на оси остаются
\[
\frac12,\quad1,\quad\sqrt2,\quad3.
\]
Точки $1$ и $\sqrt2$ выколоты.

\key{5. Метод интервалов.}
\[
t\in\left[\frac12;1\right)\cup(\sqrt2;3].
\]

\key{6. Обратная замена.}
\[
\begin{cases}
\dfrac12\le3^x<1,
\end{cases}
\qquad\text{или}\qquad
\begin{cases}
\sqrt2<3^x\le3.
\end{cases}
\]
Так как основание $3>1$,
\[
\boxed{
x\in\left[\log_3\frac12;0\right)
\cup
\left(\frac12\log_3 2;1\right].
}
\]

\section{Логарифмический переход}

Если $a^x$ сравнивается с положительным числом $b$, которое не удаётся представить степенью основания $a$, используйте логарифм:
\[
a^x\gtreqless b.
\]
При $a>1$:
\[
x\gtreqless\log_a b.
\]
При $0<a<1$ знак меняется:
\[
x\lesseqgtr\log_a b.
\]

Например,
\[
3^x>\sqrt2
\iff x>\log_3\sqrt2
\iff x>\frac12\log_3 2.
\]

Ещё один пример:
\[
5^{2x+1}\le10.
\]
Так как $5>1$, логарифмирование сохраняет знак:
\[
2x+1\le\log_5 10
\iff
\boxed{x\le\frac{\log_5 10-1}{2}}.
\]

\section{Универсальный алгоритм}

\begin{enumerate}
  \item Проверьте, что основания степеней положительны и отличны от единицы.
  \item Выпишите ограничения из знаменателей и других потенциально опасных выражений.
  \item Перепишите составные числа, дроби и корни в степенном виде.
  \item Попробуйте получить одинаковые основания в обеих частях.
  \item Если повторяется $a^x$, выполните замену $t=a^x$ и сразу запишите $t>0$.
  \item Приведите полученное неравенство к виду с нулём справа.
  \item Решите алгебраическое неравенство, сохраняя выколотые точки знаменателя.
  \item Выполните обратную замену. Каждый полученный промежуток для $t$ запишите отдельной системой.
  \item Сравните степени по монотонности основания или примените ло\-га\-риф\-ми\-че\-ский переход.
  \item Проверьте строгие границы, ограничения и объедините все подходящие промежутки.
\end{enumerate}

\section{Типичные ошибки}

\begin{itemize}
  \item \key{Отбросить разные основания.} Сравнивать показатели напрямую разрешено только при одинаковом основании.
  \item \key{Забыть сменить знак.} При $0<a<1$ показательная функция убывает.
  \item \key{Потерять условие \(t>0\).} Оно исключает лишнюю часть решения алгебраического неравенства.
  \item \key{Раскрыть знаменатель без необходимости.} Факторизованный знаменатель удобнее для метода интервалов.
  \item \key{Включить нуль знаменателя.} Запрещённая точка остаётся выколотой при любом знаке неравенства.
  \item \key{Поставить скобки неверно.} При вычитании дроби весь её числитель заключается в скобки.
  \item \key{Домножить на минус и сохранить знак.} При умножении на от\-ри\-ца\-тель\-но\-е число знак меняется.
  \item \key{Смешать системы и совокупности.} Один промежуток задаётся системой условий; несколько альтернативных промежутков объединяются совокупностью.
  \item \key{Получить десятичное приближение слишком рано.} Ответы вида $\log_a b$ сохраняют точность.
\end{itemize}

\section{Итоговый чек-лист}

\begin{itemize}
  \item[$\square$] Я узнаю показательное неравенство по переменной в показателе.
  \item[$\square$] Я проверяю условия $a>0$ и $a\ne1$.
  \item[$\square$] Я определяю, возрастает или убывает показательная функция.
  \item[$\square$] Я умею представлять корни, дроби и составные числа как степени.
  \item[$\square$] Я сравниваю показатели только после получения одинакового основания.
  \item[$\square$] Я ввожу $t=a^x$ вместе с условием $t>0$.
  \item[$\square$] Я сохраняю все ограничения знаменателя.
  \item[$\square$] Я решаю рациональное неравенство методом интервалов.
  \item[$\square$] Я выполняю обратную замену отдельными системами.
  \item[$\square$] Я применяю логарифм, когда общее основание получить неудобно.
  \item[$\square$] Я проверяю границы и записываю объединение промежутков.
\end{itemize}

\clearpage
\section{Тренировочный блок}

\textit{Решите неравенства. Первые три задания имеют средний уровень, задания 4--9 --- уровень выше среднего, задание 10 --- сложный уровень.}

\begin{enumerate}
  \item $2^x\ge16$.
  \item $\left(\dfrac13\right)^{2x-1}<\dfrac1{27}$.
  \item $4^x\cdot8<\sqrt[3]{16}$.
  \item $\sqrt2\cdot8^x\ge\dfrac4{\sqrt[5]{16}}$.
  \item $\left(\dfrac15\right)^{\frac{2x+3}{x-2}}>5$.
  \item $\left(\dfrac12\right)^{\frac{x}{x+1}}\le2$.
  \item $9^x-10\cdot3^x+9\le0$.
  \item $4^x-5\cdot2^x+4>0$.
  \item $\dfrac7{9^x-2}\ge\dfrac2{3^x-1}$.
  \item $\dfrac{5^{2x}-26\cdot5^x+25}{5^x-5}\ge0$.
\end{enumerate}

\clearpage
\section{Ответы}

\begin{longtable}{>{\bfseries}p{12mm}p{0.78\linewidth}}
\toprule
№ & Ответ \\
\midrule
\endhead
1 & $[4;+\infty)$ \\
2 & $(2;+\infty)$ \\
3 & $\left(-\infty;-\dfrac56\right)$ \\
4 & $\left[\dfrac7{30};+\infty\right)$ \\
5 & $\left(-\dfrac13;2\right)$ \\
6 & $(-\infty;-1)\cup\left[-\dfrac12;+\infty\right)$ \\
7 & $[0;2]$ \\
8 & $(-\infty;0)\cup(2;+\infty)$ \\
9 & $\left[\log_3\dfrac12;0\right)\cup\left(\dfrac12\log_3 2;1\right]$ \\
10 & $[0;1)\cup[2;+\infty)$ \\
\bottomrule
\end{longtable}

\vfill
\begin{center}
\textcolor{softgray}{Код самопроверки: \textbf{3152}}
\end{center}

\end{document}